\section{The smallest possible market: a two-step tree}
\label{sec:deamtree}

We need a machine that evaluates the best-plan value
\eqref{eq:deamstop}.  What makes the problem interesting is that at
every moment the holder compares ``stop now'' --- whose value,
intrinsic, is trivial --- against ``continue'', whose value is a
conditional expectation of everything that may happen later.  The
simplest honest machine for carrying such expectations is the binomial
tree of Cox, Ross and Rubinstein \citep{CoxRossRubinstein1979}, and it
is small enough that this section builds one from scratch and solves
it by hand.

Chop the calendar time to expiry into $N_t$ steps of length $\Delta
t=t/N_t$.  Over one step, let the spot do one of exactly two things:
multiply by $e^{\sigma\sqrt{\Delta t}}$ (an ``up'' move) or by its
reciprocal $e^{-\sigma\sqrt{\Delta t}}$ (a ``down'' move).  Why this
particular size?  Because each step then moves the \emph{log} of the
spot by $\pm\sigma\sqrt{\Delta t}$, contributing $\sigma^2\Delta t$ of
squared log-move; over $N_t$ independent steps the contributions add
to $\sigma^2 t$ --- the total variance a volatility $\sigma$ is
supposed to generate over $t$ years of calendar time.  One matching
condition down, one to go: the up-move probability $p$ is fixed by
requiring the tree to grow the spot at the carry --- the rate $r-q_d$
that links spot to forward, with $q_d$ the dividend yield exactly as
in Chapter~6:
\begin{equation}
p\,e^{\sigma\sqrt{\Delta t}}+(1-p)\,e^{-\sigma\sqrt{\Delta t}}
=e^{(r-q_d)\Delta t}
\qquad\Longrightarrow\qquad
p=\frac{e^{(r-q_d)\Delta t}-e^{-\sigma\sqrt{\Delta t}}}
       {e^{\sigma\sqrt{\Delta t}}-e^{-\sigma\sqrt{\Delta t}}}\,.
\label{eq:deamprob}
\end{equation}
For $p$ to be a probability we need $0<p<1$, which by
\eqref{eq:deamprob} says the carry factor must land strictly between
the down and up factors:
\[
e^{-\sigma\sqrt{\Delta t}}<e^{(r-q_d)\Delta t}<e^{\sigma\sqrt{\Delta t}}
\quad\Longleftrightarrow\quad
|r-q_d|\,\Delta t<\sigma\sqrt{\Delta t}
\quad\Longleftrightarrow\quad
\sigma>|r-q_d|\sqrt{t/N_t}\,.
\]
The one-step move must be able to outrun the drift.  Plug in working
numbers --- a carry gap of $4\%$, half a year, and $N_t=256$, the
depth the chapter uses when it inverts whole chains
(appendix~\ref{app:deamprotocol}) --- and the floor is
$0.04\sqrt{0.5/256}\approx0.18\%$ of volatility: the condition
excludes nothing any equity board quotes, but the inversion of
\cref{sec:deamsub} keeps its brackets above it so that every tree it
prices is a valid one.

Two structural facts make the lattice cheap.  Multiplication commutes,
so up-then-down and down-then-up land on the same spot: after $n$
steps of which $j$ were up, the spot is
$S\,e^{(2j-n)\sigma\sqrt{\Delta t}}$, and layer $n$ holds only $n+1$
distinct values (the tree \emph{recombines}).  And the value can be
computed backward.  At expiry the option is worth its payoff at each
terminal node.  One layer earlier, the holder at node $j$ compares
stopping against continuing, and continuing is worth the discounted
$p$-average of the two successor values.  Writing $A^{(n)}_j$ for the
American value at layer $n$, node $j$:
\begin{equation}
A^{(N_t)}_j=\text{payoff}_j,
\qquad
A^{(n)}_j=\max\Big\{\text{intrinsic}^{(n)}_j,\;
e^{-r\Delta t}\big(p\,A^{(n+1)}_{j+1}+(1-p)\,A^{(n+1)}_{j}\big)\Big\}.
\label{eq:deamroll}
\end{equation}
This is dynamic programming in its plainest form: start where the
answer is known and walk backward, letting each node pick the better
of its two options.  Delete the $\max$ against intrinsic and the same
recursion prices the European twin $E^{(n)}_j$ --- the two legs of
\cref{sec:deamright}'s decomposition come from one lattice, differing
by a single comparison.  Readers of Chapter~4 will recognize the
family: each value is a nonnegative combination of the next layer's,
and $\max$ preserves order, so the rollback is a monotone scheme ---
raise any input and no output falls --- which is the discrete reason
tree prices are stable and free of invented oscillation.

Now solve one completely.  Take an American put with $S=100$, $K=105$,
$\sigma=30\%$, half a year to expiry, $r=6\%$, no dividends, and just
$N_t=2$ steps --- the smallest tree in which early exercise can
happen strictly inside the lattice.  The step is $\Delta t=0.25$.  The
up factor is $e^{0.30\sqrt{0.25}}=e^{0.15}=\MacDeamTreeUpFactor$; the
down factor is its reciprocal $\MacDeamTreeDownFactor$.  The one-step
carry factor is $e^{0.06\times0.25}=\MacDeamTreeGrowth$ and the
one-step discount is $e^{-0.015}=\MacDeamTreeStepDisc$.  The
up-probability \eqref{eq:deamprob} comes out
\[
p=\frac{\MacDeamTreeGrowth-\MacDeamTreeDownFactor}
       {\MacDeamTreeUpFactor-\MacDeamTreeDownFactor}
 =\MacDeamTreeProbUp .
\]
The spot lattice is: after one step, $\MacDeamTreeSpotUp$ or
$\MacDeamTreeSpotDown$; after two, $\MacDeamTreeSpotUU$, back to
$100.00$, or $\MacDeamTreeSpotDD$.  The terminal payoffs of the
$K=105$ put are $0$, $\MacDeamTreePayUD$ and $\MacDeamTreePayDD$.
Follow \cref{fig:deamtree}(a) node by node.  At the up node
(spot $\MacDeamTreeSpotUp$), continuing is worth
\[
\MacDeamTreeStepDisc\times
\big(\MacDeamTreeProbUp\times0
+\MacDeamTreeProbDown\times\MacDeamTreePayUD\big)=\MacDeamTreeContUp ,
\]
and the put is out of the money there (intrinsic $0$), so the node's
value is $\MacDeamTreeContUp$.  At the down node (spot
$\MacDeamTreeSpotDown$) --- the orange node, and the whole lesson ---
continuing is worth
\[
\MacDeamTreeStepDisc\times\big(\MacDeamTreeProbUp\times\MacDeamTreePayUD
+\MacDeamTreeProbDown\times\MacDeamTreePayDD\big)=\MacDeamTreeContDown ,
\]
but exercising pays $105-\MacDeamTreeSpotDown=\MacDeamTreeIntrDown$
immediately.  The holder exercises: the node is worth
$\MacDeamTreeIntrDown$, and the best plan has stopped strictly before
expiry.  Rolling back once more, the root is worth
\[
\MacDeamTreeStepDisc\times\big(\MacDeamTreeProbUp\times\MacDeamTreeContUp
+\MacDeamTreeProbDown\times\MacDeamTreeIntrDown\big)=\MacDeamTreeRootAm ,
\]
comfortably above the root's own intrinsic of $5$.  Now delete the
$\max$ and reprice: the down node keeps its continuation value
$\MacDeamTreeValDownEu$, and the root comes out
$\MacDeamTreeRootEu$.  The difference,
\[
A-E=\MacDeamTreeRootAm-\MacDeamTreeRootEu=\MacDeamTreePremTwoStep ,
\]
is the first early-exercise premium this book has computed by hand ---
earned entirely at one node, where interest on the strike beat the
remaining time value.  Every number above is
reproducible with a pocket calculator, and checking one or two is the
best way to make the rollback \eqref{eq:deamroll} feel inevitable.

\begin{figure}[htbp]
\centering
\paperfig{\textwidth}{fig_deam_tree.pdf}
\caption{(a)~The hand-solved two-step American put ($S=100$, $K=105$,
$\sigma=30\%$, $t=0.5$, $r=6\%$): spot above each node, option value
inside it.  At the orange node intrinsic $\MacDeamTreeIntrDown$ beats
continuation $\MacDeamTreeContDown$ and the best plan stops early;
the European twin keeps $\MacDeamTreeValDownEu$ there, and the root
difference is the premium $\MacDeamTreePremTwoStep$.  (b)~Convergence
of the same put's price in the step count (even $N_t$, against an
$N_t=4001$ reference): both legs converge like $1/N_t$ (the dotted
guide) with an oscillatory fine structure, and the \emph{difference}
$A-E$ (green) runs below either leg --- in the median across depths
beyond $64$, \MacDeamConvRatio$\times$ smaller.}
\label{fig:deamtree}
\end{figure}

Two steps prove the mechanism; accuracy needs more.  Panel~(b) of
\cref{fig:deamtree} reprices the same put at even step counts from $2$
to $400$ and plots the absolute error against a deep reference tree.
Read three things off it.  Both the American and the European leg
converge like $1/N_t$ --- halve the step, halve the error --- with an
oscillatory fine structure that is a known cosmetic feature of
binomial lattices (the strike weaves between node levels as $N_t$
changes).  At the whole-chain working depth $N_t=256$ the American
leg is within \MacDeamConvAmErrCents\ cents of converged.  And the green
curve --- the error of the \emph{difference} $A-E$ --- runs below
either leg's, in the median \MacDeamConvRatio$\times$ smaller: the
two legs share one lattice, so most of their discretization error is
common and cancels in the subtraction.  Keep that asymmetry in mind;
it is the quiet reason the next section's method reports cleaner
numbers than a $1/N_t$ scheme has any right to.

We can now price the stopping right to known accuracy.  The next
question is how to remove it from a quote we did not price ourselves.
